\section{Basic differentiation rules}

The derivative is defined by a limit, but we do not want to compute that limit from the beginning every time. Differentiation rules allow us to compute derivatives efficiently. These rules should be read as shortcuts justified by the limit definition, not as unrelated formulas.

\subsection*{Constant functions}

A constant function does not change. Therefore its derivative is zero.

\begin{theorembox}
If
\[
f(x)=c,
\]
where \(c\) is constant, then
\[
f'(x)=0.
\]
\end{theorembox}

\begin{examplebox}
If
\[
f(x)=7,
\]
then
\[
f'(x)=0.
\]
The graph is a horizontal line, and its slope is zero everywhere.
\end{examplebox}

\subsection*{The power rule}

One of the most important rules is the power rule.

\begin{theorembox}
For a real power \(p\), where the expression is defined,
\[
\frac{d}{dx}x^p=p x^{p-1}.
\]
\end{theorembox}

For positive integers, this rule follows from expanding the difference quotient. In practice, it is used for many powers, including negative and fractional powers, on appropriate domains.

\begin{examplebox}
Differentiate
\[
f(x)=x^5.
\]
By the power rule,
\[
f'(x)=5x^4.
\]
\end{examplebox}

\begin{examplebox}
Differentiate
\[
g(x)=\frac{1}{x^3}.
\]
First rewrite
\[
g(x)=x^{-3}.
\]
Then
\[
g'(x)=-3x^{-4}=-\frac{3}{x^4}.
\]
\end{examplebox}

\begin{examplebox}
Differentiate
\[
h(x)=\sqrt{x}.
\]
Rewrite
\[
h(x)=x^{1/2}.
\]
Then
\[
h'(x)=\frac12x^{-1/2}=\frac{1}{2\sqrt{x}},
\]
for \(x>0\).
\end{examplebox}

The rewriting step is often the main work. Roots and fractions become easier to differentiate when written as powers.

\subsection*{Constant multiple rule}

Constants can be pulled out of differentiation.

\begin{theorembox}
If \(c\) is constant, then
\[
\frac{d}{dx}\bigl(c f(x)\bigr)=c f'(x).
\]
\end{theorembox}

\begin{examplebox}
If
\[
f(x)=5x^7,
\]
then
\[
f'(x)=5\cdot7x^6=35x^6.
\]
\end{examplebox}

The constant scales the function, so it also scales the rate of change.

\subsection*{Sum and difference rules}

The derivative of a sum is the sum of the derivatives.

\begin{theorembox}
If \(f\) and \(g\) are differentiable, then
\[
(f+g)'=f'+g'
\]
and
\[
(f-g)'=f'-g'.
\]
\end{theorembox}

\begin{examplebox}
Differentiate
\[
y=5x^{15}-x^2+\frac13x-2.
\]
Differentiate term by term:
\[
y'=75x^{14}-2x+\frac13.
\]
The derivative of the constant \(-2\) is zero.
\end{examplebox}

This example is typical for polynomial expressions. We read the function term by term, apply the power rule, and combine the results.

\subsection*{A useful habit}

Before differentiating, rewrite the expression into a form that matches the rules. Fractions and roots often become powers.

\begin{examplebox}
Differentiate
\[
y=\frac{2}{\sqrt{x}}-\frac{3}{\sqrt{x^3}}.
\]
Rewrite:
\[
y=2x^{-1/2}-3x^{-3/2}.
\]
Now apply the power rule:
\[
y'=2\left(-\frac12\right)x^{-3/2}-3\left(-\frac32\right)x^{-5/2}.
\]
Thus
\[
y'=-x^{-3/2}+\frac92x^{-5/2}.
\]
\end{examplebox}

\begin{summarybox}
For basic differentiation, ask:
\begin{itemize}
\item Can roots or fractions be rewritten as powers?
\item Is there a constant factor?
\item Can I differentiate term by term?
\item Which rule applies to each term?
\item On what domain does the derivative formula make sense?
\end{itemize}
\end{summarybox}

Differentiation rules are efficient because they preserve the meaning of the derivative while avoiding repeated limit computations.